\section{Method: an evaluation ladder with explicit scope}
\label{sec:method}

\subsection{The scope discipline}

Every measurement carries one of seven scopes, enforced in code rather than promised in
prose. Four matter for the results below: \texttt{deployable} (end-to-end performance of
a shippable path; no oracle, no human key in the answer path, and the denominator is not
oracle-selected); \texttt{oracle\_free\_component\_eval} (the prediction path is
oracle-free, but the \textbf{denominator is oracle-selected} --- a component result
only); \texttt{delivered} (the pipeline's own output, scored against a key); and
\texttt{identity\_oracle} (consumes human-supplied identity --- a \textbf{bound}, never
system performance). The remaining three --- \texttt{proposal\_ceiling},
\texttt{definition-change} and \texttt{bug-diagnostic} --- are defined in the
supplement, along with the full row counts.

Answer production and grading are separate functions: no layer that produces an answer
may read an expected value, and a syntax-tree test fails the build if a deployable answer
function so much as binds a key-shaped name.

\subsection{The ladder}

Five layers answer the same questions over the same stored relations, differing only in
where object identity comes from: a geometry ceiling and a stored-edge replay, both
consuming human identity; the delivered graph, using its own learned labels; a grounded
graph, using an oracle-free bridge; and direct multiview RGB, which uses no graph at all.
\Cref{fig:ladder} shows which pipeline stages each layer relies on.

\begin{figure}[t]
  \centering
  \includegraphics[width=\linewidth]{figures/fig1_evaluation_ladder.pdf}
  \caption{The evaluation ladder. Each layer holds the pipeline fixed and changes one
  source of identity; bars span the stages a layer relies on. Hatched layers consume
  human-supplied identity and are bounds, never system performance.}
  \label{fig:ladder}
\end{figure}

\subsection{Predeclaration}

Gates, stopping rules and interpretation limits are written and committed before the
corresponding scores exist. Twice this caught a \emph{structurally unreachable} outcome
before measurement: a predeclared clause that cannot fire is a broken instrument, not a
null result, and only a check in advance reveals it. In the first relation experiment the
routing structure made the accuracy clause unreachable, and it duly came in at exactly
$0.0000$ \rid{F25} with the proceed decision \texttt{false} \rid{F27}.

Blinded responses are generated in an isolated context with \textbf{no access to the key}
and hash-pinned before scoring. The protection is isolation plus the hash pin, not
version-history ordering: for the transfer run the key was recorded before the response
was pinned, so commit order does not prove the response predates the key. The full
chronology and the transfer protocol's amendment history are in the supplement.
